\documentclass[11pt,a4paper]{article}
\usepackage[margin=1in]{geometry}
\usepackage{booktabs}
\usepackage{enumitem}
\usepackage{hyperref}
\usepackage{xcolor}
\usepackage{tikz}
\usepackage{amsmath}
\usepackage{longtable}
\usepackage{graphicx}

\definecolor{best}{RGB}{0,128,0}
\definecolor{warn}{RGB}{200,100,0}
\definecolor{bad}{RGB}{200,0,0}

\title{\textbf{MolSSD Cloud Training Plan}\\[0.5em]
\large Pre-training on OMol25 \& Transfer Learning to UV Absorbers}
\author{MolSSD Project}
\date{March 2026}

\begin{document}
\maketitle
\tableofcontents
\newpage

% ============================================================================
\section{Executive Summary}
% ============================================================================

This document evaluates strategies for scaling MolSSD training beyond QM9/GEOM-Drugs
to the OMol25 dataset (83M molecules) and transitioning to UV-absorber discovery.
We compare four training strategies, provide cloud cost estimates, and give a
step-by-step execution plan with all code ready to run.

\textbf{Recommended strategy:} \textbf{Strategy C} (Hybrid) --- Pre-train on the
OMol25 Neutral organic subset ($\sim$10--20M molecules), fine-tune on a curated
UV-absorber dataset ($\sim$50K molecules with TD-DFT spectra), then run conditional
generation campaigns.

\textbf{Key insight:} OMol25 is ground-state DFT, so it lacks excited-state data
(TD-DFT spectra, oscillator strengths). However, it \emph{does} contain the
\textbf{HOMO-LUMO gap} (a direct proxy for $\lambda_\text{max}$ via $E = hc/\lambda$),
\textbf{orbital energies}, and \textbf{partial charges} (Mulliken, L\"owdin, NBO)
--- all of which are UV-relevant properties that can be used for pre-training
conditional generation. Full TD-DFT spectra are still needed for Phase~B
fine-tuning.

% ============================================================================
\section{Strategy Evaluation}
% ============================================================================

\subsection{Strategy A: Full OMol25 Pre-training}

\begin{itemize}[nosep]
    \item \textbf{Data:} Full OMol25 (83M molecules, 2--350 atoms, 83 elements)
    \item \textbf{Training:} 500K--1M steps on 8--32 GPUs, 1--2 weeks
    \item \textbf{Then:} Fine-tune on UV-absorber dataset
\end{itemize}

\begin{tabular}{@{}lp{10cm}@{}}
\toprule
\textcolor{best}{Pros} & Most general foundation model; strongest paper claims
(``pre-trained on 83M molecules''); best transfer to any downstream task \\
\textcolor{bad}{Cons} & Enormous compute cost (\$5K--15K); 500TB download;
includes metals/salts irrelevant to UV absorbers; needs 8--32 GPUs \\
\bottomrule
\end{tabular}

\medskip
\textbf{Verdict:} Ideal for the Nature paper but expensive. Worth it if grant
funding is available.

\subsection{Strategy B: Filter OMol25 for UV-Absorber-Related Molecules}

\begin{itemize}[nosep]
    \item \textbf{Data:} Filter OMol25 for conjugated/aromatic organic molecules
    \item \textbf{Filters:} Charge=0, spin=1, elements $\in$ \{H,C,N,O,F,S,P,Cl,Br\},
          contains aromatic ring, MW 150--800
    \item \textbf{Expected size:} $\sim$5--15M molecules after filtering
\end{itemize}

\begin{tabular}{@{}lp{10cm}@{}}
\toprule
\textcolor{best}{Pros} & Focused on relevant chemistry; 5--10$\times$ smaller
than full OMol25; cheaper (\$500--2K); faster to train \\
\textcolor{bad}{Cons} & Still no UV spectra in the data (ground-state DFT only);
filtering requires downloading a significant portion of OMol25; aromatic
filter may miss novel chromophore scaffolds \\
\bottomrule
\end{tabular}

\medskip
\textbf{Verdict:} Good balance but still requires OMol25 access and
substantial filtering compute.

\subsection{Strategy C: Hybrid (Recommended)}

\begin{itemize}[nosep]
    \item \textbf{Phase 1 --- Foundation:} Pre-train on OMol25 ``Neutral'' subset
          ($\sim$4M or full neutral split) or OMol25 4M random subset
    \item \textbf{Phase 2 --- Domain:} Fine-tune on GEOM-Drugs (304K drug-like molecules)
    \item \textbf{Phase 3 --- Task:} Fine-tune on curated UV-absorber dataset
          (50K+ molecules with TD-DFT spectra for conditional generation)
\end{itemize}

\begin{tabular}{@{}lp{10cm}@{}}
\toprule
\textcolor{best}{Pros} & Uses the readily available OMol25 4M subset (manageable
download); progressive specialization; each phase is independently useful;
cheapest path to strong results (\$200--1K) \\
\textcolor{warn}{Cons} & 4M subset is smaller than full OMol25 (weaker scaling
claims); requires curating UV-absorber dataset separately \\
\bottomrule
\end{tabular}

\medskip
\textbf{Verdict: \textcolor{best}{Best strategy.}} The 3-phase progressive
training is standard practice (cf.\ GPT pre-train $\to$ fine-tune $\to$ RLHF).
The 4M subset is large enough to demonstrate scaling. The UV-absorber dataset
is needed regardless of which strategy we choose.

\subsection{Strategy D: Direct UV-Absorber Training (No Pre-training)}

\begin{itemize}[nosep]
    \item \textbf{Data:} Curate UV-absorber dataset from scratch (ChEMBL, PubChem,
          ZINC + TD-DFT calculations)
    \item \textbf{Training:} Train MolSSD directly on 50K UV-absorber molecules
\end{itemize}

\begin{tabular}{@{}lp{10cm}@{}}
\toprule
\textcolor{best}{Pros} & Cheapest (\$50--200); fastest to results; focused \\
\textcolor{bad}{Cons} & Small dataset = poor generalization; weaker paper
(no pre-training story); limited chemical diversity \\
\bottomrule
\end{tabular}

\medskip
\textbf{Verdict:} Backup plan if cloud budget is very limited. Not recommended
for a Nature paper.

% ============================================================================
\section{Strategy Comparison Table}
% ============================================================================

\begin{table}[h]
\centering
\small
\begin{tabular}{@{}lcccccc@{}}
\toprule
& \textbf{Pre-train} & \textbf{Fine-tune} & \textbf{UV Data} &
\textbf{GPUs} & \textbf{Time} & \textbf{Cost} \\
\midrule
A: Full OMol25 & 83M mols & UV-absorbers & Separate & 8--32 & 1--2 wk & \$5--15K \\
B: Filtered OMol25 & 5--15M mols & UV-absorbers & Separate & 4--8 & 3--5 d & \$500--2K \\
\textbf{C: Hybrid} & \textbf{4M mols} & \textbf{GEOM $\to$ UV} &
\textbf{Separate} & \textbf{1--4} & \textbf{2--4 d} & \textbf{\$200--1K} \\
D: Direct UV & --- & --- & 50K mols & 1 & 1--2 d & \$50--200 \\
\bottomrule
\end{tabular}
\caption{Strategy comparison. Strategy C provides the best cost/impact tradeoff.}
\end{table}

% ============================================================================
\section{Cloud Provider Analysis}
% ============================================================================

\subsection{GPU Options for MolSSD Training}

\begin{table}[h]
\centering
\small
\begin{tabular}{@{}llcccc@{}}
\toprule
\textbf{Provider} & \textbf{GPU} & \textbf{VRAM} & \textbf{\$/hr (1 GPU)} &
\textbf{\$/hr (8 GPU)} & \textbf{Notes} \\
\midrule
Lambda Labs & A100 80GB & 80GB & \$1.10 & \$8.80 & Best value, often sold out \\
Lambda Labs & H100 80GB & 80GB & \$2.49 & \$19.92 & Fastest option \\
RunPod & A100 80GB & 80GB & \$1.64 & \$13.12 & On-demand, good availability \\
RunPod & A6000 48GB & 48GB & \$0.76 & \$6.08 & Budget option, adequate \\
Vast.ai & A100 80GB & 80GB & \$0.90 & \$7.20 & Cheapest, variable quality \\
AWS & p4d.24xlarge & 8$\times$A100 & --- & \$32.77 & Reliable, expensive \\
GCP & a2-highgpu-8g & 8$\times$A100 & --- & \$29.39 & Good for academic credits \\
\bottomrule
\end{tabular}
\caption{Cloud GPU pricing (March 2026, approximate on-demand rates).}
\end{table}

\textbf{Recommendation:} Lambda Labs or RunPod with A100 80GB. For the Hybrid
strategy (Strategy C), a single A100 is sufficient for Phase 1 (OMol25 4M,
$\sim$48h) and Phase 2 (GEOM-Drugs, $\sim$24h). Use 4$\times$A100 only if
training the full OMol25.

% ============================================================================
\section{Step-by-Step Execution Plan (Strategy C)}
% ============================================================================

\subsection{Phase 0: Preparation (Local, Before Cloud)}

\begin{enumerate}[nosep]
    \item \textbf{Apply for OMol25 access} via Globus
          (\url{https://fair-chem.github.io/molecules/datasets/omol25.html}).
          Allow 1--2 weeks for approval.
    \item \textbf{Download OMol25 4M subset} ($\sim$50--100GB) via Globus transfer.
          Path: \texttt{alcf\#dtn\_eagle} $\to$ \texttt{/OMol25/train\_4M/}
    \item \textbf{Download GEOM-Drugs} ($\sim$10GB) from Harvard Dataverse.
    \item \textbf{Curate UV-absorber dataset} (see Section~\ref{sec:uv-dataset}).
    \item \textbf{Verify code locally:}
\begin{verbatim}
python3 -c "from molssd.data.omol25_loader import OMol25MolSSD; print('OK')"
python3 -c "from molssd.data.geom_drugs_loader import GEOMDrugsMolSSD; print('OK')"
\end{verbatim}
    \item \textbf{Package code} for cloud upload:
\begin{verbatim}
tar czf molssd_code.tar.gz molssd/ scripts/ CLAUDE.md
\end{verbatim}
\end{enumerate}

\subsection{Phase 1: Pre-train on OMol25 4M (Cloud, 1--2 days)}

\begin{enumerate}[nosep]
    \item \textbf{Provision:} 1$\times$A100 80GB on Lambda Labs or RunPod
    \item \textbf{Setup:}
\begin{verbatim}
# Install dependencies
pip install torch torchvision torch-geometric rdkit fairchem-core wandb

# Upload code and data
scp molssd_code.tar.gz cloud:~/
scp -r omol25_4M/ cloud:~/data/omol25/
\end{verbatim}
    \item \textbf{Train:}
\begin{verbatim}
python3 scripts/train_omol25.py \
    --lmdb-path ~/data/omol25/train_4M/ \
    --batch-size 64 \
    --max-steps 500000 \
    --max-atoms 200 \
    --lr 3e-4 \
    --checkpoint-dir ./checkpoints_omol25 \
    --wandb
\end{verbatim}
    \item \textbf{Expected:} $\sim$48h on 1$\times$A100, cost $\sim$\$50--120
    \item \textbf{Output:} \texttt{checkpoints\_omol25/checkpoint\_final.pt}
    \item \textbf{Download checkpoint} to local machine
\end{enumerate}

\subsection{Phase 2: Fine-tune on GEOM-Drugs (Cloud or Local, 1 day)}

\begin{enumerate}[nosep]
    \item \textbf{Transfer learn} from OMol25 checkpoint:
\begin{verbatim}
python3 scripts/train_geom_drugs.py \
    --pretrained checkpoints_omol25/checkpoint_final.pt \
    --batch-size 32 \
    --max-steps 200000 \
    --lr 1e-4 \
    --checkpoint-dir ./checkpoints_geom
\end{verbatim}
    \item \textbf{Expected:} $\sim$24h on 1$\times$A100 or RTX 4090
    \item \textbf{Evaluate:} Run benchmark comparison vs EDM on GEOM-Drugs
\begin{verbatim}
python3 scripts/run_evaluation.py \
    --checkpoint checkpoints_geom/checkpoint_final.pt \
    --num-molecules 10000 \
    --output-dir evaluation_geom_drugs/
\end{verbatim}
    \item \textbf{Key metric:} Demonstrate 3--4$\times$ speedup over EDM
          on larger molecules (log-log scaling plot)
\end{enumerate}

\subsection{Phase 3: Fine-tune on UV-Absorber Dataset (Cloud or Local, 1 day)}

\begin{enumerate}[nosep]
    \item \textbf{Load GEOM-Drugs checkpoint} and fine-tune with conditional generation:
\begin{verbatim}
python3 scripts/train_uv_conditional.py \
    --pretrained checkpoints_geom/checkpoint_final.pt \
    --uv-dataset ./data/uv_absorbers/ \
    --condition-on spectrum,photostability,logP \
    --classifier-free-guidance --guidance-dropout 0.15 \
    --batch-size 32 --max-steps 100000 --lr 5e-5
\end{verbatim}
    \item \textbf{Expected:} $\sim$12--24h
    \item \textbf{Generate UV-absorber candidates:}
\begin{verbatim}
python3 scripts/generate_uv_absorbers.py \
    --checkpoint checkpoints_uv/checkpoint_final.pt \
    --target-lambda-max 320 \
    --target-epsilon 30000 \
    --guidance-weight 2.5 \
    --num-molecules 10000
\end{verbatim}
\end{enumerate}

% ============================================================================
\section{UV-Absorber Dataset Curation}
\label{sec:uv-dataset}
% ============================================================================

Since OMol25 contains \textbf{no excited-state or UV absorption data}, we must
curate a separate UV-absorber dataset. This is required regardless of which
pre-training strategy we choose.

\subsection{Data Sources}

\begin{enumerate}[nosep]
    \item \textbf{ChEMBL/PubChem:} Search for compounds with UV-Vis assay data.
          Keywords: ``UV absorber'', ``sunscreen'', ``photoprotection'',
          ``benzotriazole'', ``benzophenone'', ``triazine''.
    \item \textbf{ZINC:} Commercially available UV-absorbing compounds.
    \item \textbf{Literature:} Manual curation from review articles on
          UV-absorber design (Shaath 2010, Serpone 2007).
    \item \textbf{Patent databases:} PatentLens/Google Patents for patented
          UV-absorber structures.
\end{enumerate}

\textbf{Target:} 50,000+ molecules across 8 chromophore classes:
benzotriazoles, benzophenones, triazines, cyanoacrylates, oxanilides,
cinnamates, dibenzoylmethanes, phenyl salicylates.

\subsection{TD-DFT Spectrum Computation}

For molecules without experimental UV-Vis spectra:

\begin{enumerate}[nosep]
    \item \textbf{Geometry optimization:} eSEN (from OMol25) or
          $\omega$B97M-V/def2-SVP
    \item \textbf{Excited states:} TD-DFT with CAM-B3LYP/6-311+G(d,p),
          PCM solvent model (ethanol)
    \item \textbf{Compute first 20 excited states} $\to$ extract
          $\lambda_\text{max}$, oscillator strengths, full spectrum
    \item \textbf{Batch compute} on HPC cluster (ORCA or Gaussian)
    \item \textbf{Validation:} Compare TD-DFT $\lambda_\text{max}$ vs
          experimental for 50+ known UV absorbers (target MAE $<$ 20nm)
\end{enumerate}

\subsection{UV-Absorber Dataset Format}

Each molecule in the dataset should contain:
\begin{itemize}[nosep]
    \item 3D coordinates (eSEN-optimized geometry)
    \item Atom types and bonds
    \item $\lambda_\text{max}$ (nm) --- primary conditioning target
    \item Molar extinction coefficient $\varepsilon_\text{max}$
    \item Discretized absorption spectrum (200--800nm, 1nm bins)
    \item Photostability score (binary or continuous)
    \item logP, molecular weight, SA score
    \item Chromophore class label
\end{itemize}

% ============================================================================
\section{Cost Estimates (Strategy C)}
% ============================================================================

\begin{table}[h]
\centering
\begin{tabular}{@{}lcccc@{}}
\toprule
\textbf{Phase} & \textbf{GPU} & \textbf{Hours} & \textbf{\$/hr} & \textbf{Cost} \\
\midrule
Phase 1: OMol25 4M pre-train & 1$\times$A100 & 48 & \$1.10 & \$53 \\
Phase 2: GEOM-Drugs fine-tune & 1$\times$A100 & 24 & \$1.10 & \$26 \\
Phase 3: UV-absorber fine-tune & 1$\times$A100 & 24 & \$1.10 & \$26 \\
UV-absorber generation campaigns & 1$\times$A100 & 8 & \$1.10 & \$9 \\
\midrule
\textbf{Total (Lambda Labs)} & & \textbf{104} & & \textbf{\$114} \\
\midrule
\textbf{Total (RunPod A100)} & & 104 & \$1.64 & \$170 \\
\textbf{Total (AWS p4d)} & & 104 & \$4.10$^\dagger$ & \$427 \\
\bottomrule
\multicolumn{5}{l}{\small $^\dagger$AWS per-GPU rate from p4d.24xlarge / 8} \\
\end{tabular}
\caption{Cost estimates for Strategy C (Hybrid). Total: \$114--427 depending on provider.}
\end{table}

\textbf{With academic cloud credits:} Many cloud providers offer free credits
for research. AWS (\$5K via AWS Research Credits), GCP (\$5K via Research
Credits Program), Azure (\$5K via AI for Good). These would cover the full
training multiple times over.

% ============================================================================
\section{OMol25 Properties Relevant to UV Absorption}
\label{sec:omol25-uv-properties}
% ============================================================================

While OMol25 lacks excited-state (TD-DFT) data, it contains several
ground-state DFT properties that are \textbf{directly predictive} of
UV absorption behavior. These can be used as conditioning signals during
pre-training.

\subsection{Directly Available Properties}

\begin{table}[h]
\centering
\small
\begin{tabular}{@{}llp{6cm}@{}}
\toprule
\textbf{Property} & \textbf{Source} & \textbf{UV Absorption Relevance} \\
\midrule
HOMO--LUMO gap & DFT ($\omega$B97M-V) &
    \textbf{Primary proxy for $\lambda_\text{max}$}.
    $E_\text{gap} = hc/\lambda$ gives approximate absorption wavelength.
    Gap of 3.1\,eV $\approx$ 400\,nm (UVA), 4.4\,eV $\approx$ 280\,nm (UVB). \\
HOMO energy & DFT & Electron donor capacity. Low HOMO $\to$ harder to
    excite. Relevant for identifying electron-rich chromophores. \\
LUMO energy & DFT & Electron acceptor capacity. Push-pull systems with
    low LUMO absorb at longer wavelengths. \\
Mulliken charges & DFT & Indicates intramolecular charge transfer (ICT).
    Large charge separation $\to$ strong $\pi \to \pi^*$ transitions. \\
L\"owdin charges & DFT & More basis-set-independent than Mulliken.
    Better for comparing across molecule sizes. \\
NBO charges & DFT (via GBW) & Natural bond orbital analysis. Best
    for identifying donor--acceptor groups and $n \to \pi^*$ transitions. \\
Orbital energies & DFT & Full MO spectrum identifies gaps between
    occupied/virtual orbitals for $n \to \pi^*$, $\pi \to \pi^*$,
    $\sigma \to \sigma^*$ transitions. \\
Total energy & DFT & Thermodynamic stability. Unstable molecules
    decompose under UV exposure. \\
Per-atom forces & DFT & Near-zero forces $\to$ equilibrium geometry.
    Important for accurate spectrum prediction. \\
\bottomrule
\end{tabular}
\caption{OMol25 properties relevant to UV absorption prediction.}
\end{table}

\subsection{Available via Post-Processing of ORCA Files}

\begin{table}[h]
\centering
\small
\begin{tabular}{@{}llp{6cm}@{}}
\toprule
\textbf{Property} & \textbf{Source} & \textbf{UV Relevance} \\
\midrule
Density matrix & \texttt{density\_mat.npz} &
    Encodes $\pi$-delocalization extent. Larger conjugated systems $\to$
    smaller gap $\to$ longer $\lambda_\text{max}$. \\
Fock matrix & GBW wavefunction &
    Eigenvalues give orbital energies; off-diagonal elements encode
    orbital coupling (relates to oscillator strength). \\
Molecular orbitals & GBW wavefunction &
    Spatial distribution of HOMO/LUMO determines transition dipole
    moment direction and magnitude. \\
Bond orders & NBO analysis &
    Identifies $\pi$-bond conjugation pathways critical for
    chromophore identification. \\
\bottomrule
\end{tabular}
\end{table}

\subsection{HOMO--LUMO Gap as $\lambda_\text{max}$ Proxy}

The relationship between the DFT HOMO--LUMO gap and experimental
$\lambda_\text{max}$ is:
\[
    \lambda_\text{max} \approx \frac{hc}{E_\text{HOMO-LUMO}}
    = \frac{1240}{E_\text{gap}\,(\text{eV})} \; \text{nm}
\]

\begin{table}[h]
\centering
\begin{tabular}{@{}lcc@{}}
\toprule
\textbf{UV Range} & \textbf{$\lambda$ (nm)} & \textbf{Approx.\ Gap (eV)} \\
\midrule
UV-C & 200--280 & 4.4--6.2 \\
UV-B & 280--320 & 3.9--4.4 \\
UV-A & 320--400 & 3.1--3.9 \\
Visible & 400--700 & 1.8--3.1 \\
\bottomrule
\end{tabular}
\caption{HOMO--LUMO gap ranges for target UV absorption regions.}
\end{table}

\textbf{Accuracy:} For organic chromophores, the $\omega$B97M-V HOMO--LUMO gap
correlates with TD-DFT $\lambda_\text{max}$ at $R^2 \approx 0.7$--$0.85$, with a
systematic overestimation of 0.5--1.0\,eV (underestimates $\lambda_\text{max}$ by
30--80\,nm). This is sufficient for \emph{pre-training} the conditional model;
fine-tuning with TD-DFT spectra corrects the systematic error.

\subsection{Pre-training Strategy Using OMol25 Properties}

\textbf{Phase 1 (OMol25 pre-training):} Condition MolSSD on the HOMO--LUMO gap
as a scalar property (via adaLN conditioning). This teaches the model which
molecular structures produce which gap values $\to$ approximate $\lambda_\text{max}$
control.

\textbf{Phase 3 (UV fine-tuning):} Replace HOMO--LUMO gap conditioning with
actual TD-DFT $\lambda_\text{max}$ and full spectrum conditioning. The model
transfers its learned structure--gap relationship and refines it.

\subsection{Properties NOT in OMol25 (Need Separate Computation)}

\begin{itemize}[nosep]
    \item \textbf{TD-DFT excited-state energies} --- actual $\lambda_\text{max}$
          (HOMO--LUMO gap is a proxy, not exact)
    \item \textbf{Oscillator strengths} --- determines $\varepsilon_\text{max}$
          (absorption intensity)
    \item \textbf{Full UV-Vis spectrum} --- needed for spectrum conditioning
    \item \textbf{Photostability metrics} --- excited-state lifetimes,
          conical intersections, ESIPT
    \item \textbf{Dipole moment} --- not explicitly stored (derivable from
          charges + positions)
    \item \textbf{Polarizability} --- not stored (derivable from Fock matrix
          perturbation theory)
    \item \textbf{Solvatochromic shifts} --- OMol25 is gas-phase DFT
\end{itemize}

% ============================================================================
\section{What OMol25 Pre-training Gives Us}
% ============================================================================

Pre-training on OMol25 provides:

\begin{enumerate}[nosep]
    \item \textbf{Learned 3D chemical grammar:} Bond lengths, angles, torsions
          for diverse organic chemistry
    \item \textbf{Element embeddings:} Meaningful representations for 15+
          atom types (vs 5 for QM9)
    \item \textbf{HOMO--LUMO gap awareness:} Model learns which molecular
          structures produce which electronic gaps --- a direct proxy for
          $\lambda_\text{max}$ (see Section~\ref{sec:omol25-uv-properties})
    \item \textbf{Charge-transfer intuition:} Partial charges teach the model
          about donor--acceptor systems critical for UV chromophores
    \item \textbf{Scale awareness:} Model learns to handle molecules from 2 to
          200+ atoms
    \item \textbf{Coarsening intuition:} The spectral coarsening hierarchy is
          learned across diverse molecular topologies
    \item \textbf{Faster convergence:} Fine-tuning on 50K UV absorbers converges
          in $\sim$10$\times$ fewer steps than training from scratch
    \item \textbf{Better generalization:} Pre-trained model explores chemical
          space more effectively during generation
\end{enumerate}

\textbf{Analogy:} This is like pre-training a language model on all of Wikipedia
before fine-tuning it to write poetry. The model learns grammar and vocabulary
first, then learns style.

% ============================================================================
\section{OMol25 Filtering for UV-Relevant Molecules}
% ============================================================================

If we want a smaller, more relevant pre-training subset, we can filter OMol25:

\subsection{Filter Criteria}

\begin{enumerate}[nosep]
    \item \textbf{Charge = 0, Spin = 1} (neutral singlets --- use ``Neutral'' split)
    \item \textbf{Elements $\subseteq$ \{H, C, N, O, F, S, P, Cl, Br, I\}}
          (exclude metals, lanthanides)
    \item \textbf{Molecular weight 100--800 Da} (drug-like range)
    \item \textbf{Contains $\geq 1$ aromatic ring} (UV absorption requires
          $\pi$-conjugation)
    \item \textbf{$\leq 200$ atoms} (practical limit for training)
\end{enumerate}

\subsection{Expected Yield}

\begin{itemize}[nosep]
    \item OMol25 full: 83M molecules
    \item After charge/spin filter (Neutral): $\sim$20--30M
    \item After element filter: $\sim$15--25M
    \item After MW + aromatic filter: $\sim$5--15M
    \item After atom count filter: $\sim$5--12M
\end{itemize}

This 5--12M subset is large enough for strong pre-training while being focused
on the chemistry relevant to UV absorbers.

% ============================================================================
\section{Multi-GPU Training Setup}
% ============================================================================

For training on the full OMol25 (Strategy A) or large filtered subsets:

\subsection{PyTorch DDP Configuration}

\begin{verbatim}
# Launch on 8 GPUs (single node)
torchrun --nproc_per_node=8 scripts/train_omol25.py \
    --batch-size 32 \       # per-GPU batch size (32 * 8 = 256 effective)
    --lr 1.2e-3 \           # linear scaling: 3e-4 * (256/64)
    --warmup-steps 2000 \
    --max-steps 500000 \
    --ddp

# Launch on 4 nodes x 8 GPUs (32 total)
torchrun --nnodes=4 --nproc_per_node=8 \
    --rdzv_backend=c10d --rdzv_endpoint=$MASTER_ADDR:$MASTER_PORT \
    scripts/train_omol25.py \
    --batch-size 16 \       # per-GPU (16 * 32 = 512 effective)
    --lr 2.4e-3 \           # linear scaling
    --ddp
\end{verbatim}

\subsection{Memory Requirements}

\begin{table}[h]
\centering
\begin{tabular}{@{}lcc@{}}
\toprule
\textbf{Component} & \textbf{Per GPU (A100 80GB)} & \textbf{Per GPU (4090 16GB)} \\
\midrule
Model (10M params, fp32) & 40 MB & 40 MB \\
Optimizer state (AdamW) & 120 MB & 120 MB \\
Activations + gradients & 2--8 GB & 2--4 GB \\
Batch data (bs=32, 200 atoms) & 1--2 GB & 1--2 GB \\
\textbf{Total} & \textbf{$\sim$4--10 GB} & \textbf{$\sim$3--6 GB} \\
\textbf{Available} & 80 GB & 16 GB \\
\bottomrule
\end{tabular}
\caption{Memory requirements. Even the RTX 4090 has ample headroom.}
\end{table}

% ============================================================================
\section{Timeline}
% ============================================================================

\begin{table}[h]
\centering
\begin{tabular}{@{}clll@{}}
\toprule
\textbf{Week} & \textbf{Task} & \textbf{Where} & \textbf{Deliverable} \\
\midrule
1 & Apply for OMol25 access & Local & Globus approval \\
1 & Curate UV-absorber molecules & Local & 50K SMILES + metadata \\
2 & Download OMol25 4M + GEOM-Drugs & Local & Raw data on disk \\
2 & Run TD-DFT on UV absorbers & HPC cluster & Spectra for 50K mols \\
3 & Phase 1: Pre-train on OMol25 4M & Cloud (A100) & Pre-trained checkpoint \\
3 & Phase 2: Fine-tune GEOM-Drugs & Cloud (A100) & Scaling benchmark \\
4 & Phase 3: Fine-tune UV-absorbers & Cloud (A100) & Conditional model \\
4 & Generation campaigns & Cloud (A100) & 50K+ candidates \\
5 & Screening pipeline & Local & Top 500 verified \\
5 & Paper: fill tables \& figures & Local & Draft ready \\
\bottomrule
\end{tabular}
\caption{5-week execution timeline for Strategy C.}
\end{table}

% ============================================================================
\section{Code Readiness Checklist}
% ============================================================================

\begin{table}[h]
\centering
\begin{tabular}{@{}lcc@{}}
\toprule
\textbf{Component} & \textbf{Status} & \textbf{File} \\
\midrule
QM9 data loader & \textcolor{best}{Done} & \texttt{molssd/data/qm9\_loader.py} \\
GEOM-Drugs data loader & \textcolor{best}{Done} & \texttt{molssd/data/geom\_drugs\_loader.py} \\
OMol25 data loader & \textcolor{best}{Done} & \texttt{molssd/data/omol25\_loader.py} \\
QM9 training script & \textcolor{best}{Done} & \texttt{scripts/train\_qm9.py} \\
GEOM-Drugs training script & \textcolor{best}{Done} & \texttt{scripts/train\_geom\_drugs.py} \\
OMol25 training script & \textcolor{warn}{TODO} & \texttt{scripts/train\_omol25.py} \\
UV-absorber training script & \textcolor{warn}{TODO} & \texttt{scripts/train\_uv\_conditional.py} \\
Sampling script & \textcolor{best}{Done} & \texttt{scripts/sample.py} \\
Evaluation script & \textcolor{best}{Done} & \texttt{scripts/run\_evaluation.py} \\
Ablation script & \textcolor{best}{Done} & \texttt{scripts/run\_ablations.py} \\
Spectrum conditioning module & \textcolor{bad}{TODO} & \texttt{molssd/models/conditioning.py} \\
Classifier-free guidance & \textcolor{bad}{TODO} & (in training loop) \\
DDP wrapper & \textcolor{bad}{TODO} & (in training script) \\
UV-absorber data loader & \textcolor{bad}{TODO} & \texttt{molssd/data/uv\_loader.py} \\
\bottomrule
\end{tabular}
\caption{Code readiness for cloud training. Green = ready, orange = straightforward,
red = needs design work.}
\end{table}

% ============================================================================
\section{Recommendation}
% ============================================================================

\begin{enumerate}
    \item \textbf{Finish QM9 training} (currently running, $\sim$3h remaining)
          and evaluate (Step A3).
    \item \textbf{Apply for OMol25 access now} --- approval takes 1--2 weeks,
          don't wait.
    \item \textbf{Start UV-absorber dataset curation} in parallel (ChEMBL/PubChem
          queries, literature search).
    \item \textbf{Run GEOM-Drugs locally} on your RTX 4090 (Step A5) while
          waiting for OMol25 access.
    \item \textbf{When OMol25 is available:} Spin up 1$\times$A100 on Lambda Labs,
          run Phases 1--3 in sequence ($\sim$4 days, $\sim$\$114).
    \item \textbf{Apply for cloud credits} from AWS/GCP/Azure research programs
          to cover costs.
\end{enumerate}

\end{document}
